% Part: sets-functions-relations
% Chapter: induction
% Section: induction

\documentclass[../../../include/open-logic-section]{subfiles}

\begin{document}

\olfileid{sfr}{ind}{int}

\olsection{ভূমিকা}

\begin{explain}
গণিত ও যুক্তিবিদ্যায় আরোহ একটি বহুল ব্যবহৃত প্রমাণ-কৌশল।
গণিতে স্বাভাবিক সংখ্যার উপর আরোহের সঙ্গে তোমাদের পরিচয় থাকতে পারে।
গাণিতিক ও যৌক্তিক বস্তুর আরও বহু ধরনের সেটেও আরোহ প্রয়োগ করা যায়।

সাধারণভাবে আরোহ একটি সার্বিক দাবি প্রমাণ করতে দেয়; অর্থাৎ,
একটি নির্দিষ্ট ধরনের প্রত্যেক বস্তুর একটি ধর্ম আছে বলে দেখায়।
বিশেষত, “সর্বসূচক প্রবর্তন” পদ্ধতিতে প্রমাণ অত্যন্ত জটিল হলে
আরোহ প্রায়ই কাজে লাগে। আরোহ কেবল বিশেষভাবে গঠিত গাণিতিক
বস্তুর ক্ষেত্রেই কাজ করে: সেটের প্রত্যেক উপাদান হয় মৌলিক,
নয় মৌলিক উপাদান থেকে গঠিত হলে সেই সেটে আরোহ ব্যবহার করা যায়।
আরও নির্দিষ্টভাবে:
\end{explain}

\begin{defn}
একটি সেট $S$ অপেক্ষক $f$-এর অধীনে \emph{বদ্ধ}, যদি এবং কেবল যদি
$x \in S$ হলেই $f(x) \in S$ হয়।
\end{defn}

\begin{explain}
যেমন, স্বাভাবিক সংখ্যার সেট ($\mathbb{N}$) উত্তরসূরি অপেক্ষক
$f(x) = x+1$-এর অধীনে বদ্ধ। $x$ স্বাভাবিক সংখ্যা হলে $x+1$-ও তাই।
কিন্তু স্বাভাবিক সংখ্যার সেট পূর্বসূরি অপেক্ষক $f(x) = x-1$-এর
অধীনে বদ্ধ নয়: $0$ স্বাভাবিক সংখ্যা হলেও $-1$ তা নয়।
\end{explain}

\begin{defn}
কোনো ধর্ম $P$ অপেক্ষক $f(x)$-এর অধীনে \emph{সংরক্ষিত হয়}, যদি
$f$-এর সংজ্ঞাক্ষেত্রের প্রত্যেক বস্তু $a$-র জন্য $P(a)$ থেকে
$P(f(a))$ অনুসৃত হয় ($a$-র ধর্ম $P$ থাকলে $f(a)$-রও তা থাকে)।
\end{defn}

\begin{explain}
“জোড় সংখ্যা হওয়া” ধর্মটি অপেক্ষক $f(x) = x+2$-এর অধীনে
সংরক্ষিত হয়: $x$ জোড় হলে $x+2$-ও জোড়।
কিন্তু ধর্মটি উত্তরসূরি অপেক্ষকের অধীনে সংরক্ষিত হয় না:
$x$ জোড় হলে $x+1$ বিজোড়।

স্বাভাবিক সংখ্যায় আরোহ কাজ করে, কারণ 0 থেকে উত্তরসূরি অপেক্ষক
সসীমসংখ্যক বার প্রয়োগ করে প্রত্যেক স্বাভাবিক সংখ্যায় পৌঁছানো যায়।
আরোহ প্রয়োগযোগ্য প্রত্যেক সেটেরই এই ধরনের বৈশিষ্ট্য থাকে।
\end{explain}

\begin{defn}[$\Nat$-এর উপর আরোহ]
0-র কোনো ধর্ম $P$ থাকলে এবং উত্তরসূরি অপেক্ষকের অধীনে
$P$ সংরক্ষিত হলে প্রত্যেক স্বাভাবিক সংখ্যার ধর্ম $P$ থাকে।
\end{defn}

\begin{ex} শূন্য থেকে $n$ পর্যন্ত স্বাভাবিক সংখ্যাগুলির যোগফল
$\frac{n(n+1)}{2}$। \\

\textbf{ভিত্তিক্ষেত্র।} শূন্য থেকে শূন্য পর্যন্ত যোগফল হল
$0=\frac{0\cdot 1}{2}$।\\

\textbf{আরোহী ধাপ।} ধরা যাক $k$-র জন্য ধর্মটি সিদ্ধ।
দেখাব, $k+1$-এর জন্যও তা সিদ্ধ। অর্থাৎ, ধরা যাক $P(k)$ সিদ্ধ:
\[ 0 + 1 + \ldots + k = \frac{k(k+1)}{2} \]
দেখাব যে $P(k+1)$ সিদ্ধ, অর্থাৎ
\[ 0 + 1 + \ldots + k + (k+1) = \frac{(k+1)(k+2)}{2} \]
এ জন্য $P(k)$ অনুমানটি ব্যবহার করি।

\begin{align*}
0 + 1 + \ldots + k &= \frac{k(k+1)}{2} \tag{\text{অনুমান}} \\
0 + 1 + \ldots + k + (k+1) &= \frac{k(k+1)}{2} + (k+1)
\tag{\text{উভয় পক্ষে }k+1\text{ যোগ করি}} \\
&= \frac{k(k+1) + 2(k+1)}{2} \\
&= \frac{k^2 + k + 2k +2}{2} \\
&=\frac{(k+1)(k+2)}{2}
\end{align*}
সুতরাং আরোহী ধাপটি সিদ্ধ হয়েছে, এবং দাবিটি প্রমাণিত।
\end{ex}

\begin{explain}
!!{formula}s-এর ক্ষেত্রে মৌলিক উপাদানগুলি হল পরমাণবিক
!!{formula}s। অপেক্ষাকৃত সরল !!{formula}s-কে অপারেটরের সাহায্যে
যুক্ত করে জটিলতর !!{formula}s পাওয়া যায়। প্রত্যেক অপারেটরকে
তার প্রয়োজনীয় সংখ্যক সূত্র-আর্গুমেন্ট থেকে নতুন
!!{formula} গঠনকারী অপেক্ষক হিসেবে দেখা যায়। যেমন, দুটি
!!{formula}s $!A$ ও $!B$-কে সংযোজনে যুক্ত করার অপেক্ষকটি
$\mathcal E _\land (!A, !B) = !A \land !B$ হিসেবে লেখা যায়।
এগুলিকে \emph{সূত্র-গঠনকারী অপেক্ষক} বলা যায়।
!!{formula}s-এর সেট এই অপেক্ষকগুলির অধীনে বদ্ধ:
$!A$ ও $!B$ !!{formula}s হলে $\lnot !A$,
$!A \land !B$ ইত্যাদিও !!{formula}s।
\end{explain}

\begin{defn}[!!{formula}s-এর উপর আরোহ]
প্রত্যেক পরমাণবিক !!{formula}-র ধর্ম $P$ থাকলে এবং
!!{formula}-গঠনকারী অপেক্ষকগুলির অধীনে $P$ সংরক্ষিত হলে
প্রত্যেক !!{formula}-র ধর্ম $P$ থাকে।
\end{defn}

\begin{explain}
এই সংজ্ঞা !!{formula}s-এর উপর আরোহের প্রমাণের একটি পদ্ধতি দেয়।
প্রত্যেক !!{formula}-র ধর্ম $P$ আছে বলে দেখাতে হলে নিচের
ধাপগুলি অনুসরণ করো:\\

\textbf{ভিত্তিক্ষেত্র।} ধরা যাক $!A$ একটি পরমাণবিক !!{formula}।
[\ldots] সুতরাং $!A$-র ধর্ম $P$ আছে।

\textbf{আরোহী ধাপ।} ধরা যাক $!A$ ও $!B$ দুটি !!{formula}s,
এবং উভয়েরই ধর্ম $P$ আছে।

$\lnot$-ক্ষেত্র: [\ldots] সুতরাং $\lnot !A$-র ধর্ম $P$ আছে।

$\land$-ক্ষেত্র: [\ldots] সুতরাং $!A \land !B$-র ধর্ম $P$ আছে।

$\lor$-ক্ষেত্র: [\ldots] সুতরাং $!A \lor !B$-র ধর্ম $P$ আছে।

$\lif$-ক্ষেত্র: [\ldots] সুতরাং $!A \lif !B$-র ধর্ম $P$ আছে।

$\liff$-ক্ষেত্র: [\ldots] সুতরাং $!A \liff !B$-র ধর্ম $P$ আছে।

$\lforall[]$-ক্ষেত্র: [\ldots] সুতরাং $\lforall[x] !A$-র ধর্ম $P$ আছে।

$\lexists[]$-ক্ষেত্র: [\ldots] সুতরাং $\lexists[x] !A$-র ধর্ম $P$ আছে। \\

অতএব প্রত্যেক !!{formula}-র ধর্ম $P$ আছে।
\end{explain}
% BN-SRC-678 TARGET-CORRECTION-BEGIN
% Actual include-depth repair retained; implementation note hidden.
% BN-SRC-678 TARGET-CORRECTION-END
% BN-SRC-679 TARGET-CORRECTION-BEGIN
\par\smallskip\noindent\textnormal{\small\textbf{উৎস-সংশোধন (BN-SRC-679):} যোগফলের বীজগাণিতিক বিস্তারে উৎসের প্রথম 2k পদটি আসলে k-এর বর্গ; সংশোধনে পরবর্তী উৎপাদকীকরণটি সঠিক হয়।} % BN-SRC-679 TARGET-CORRECTION-END
% BN-SRC-680 TARGET-CORRECTION-BEGIN
% Explicit summation endpoints follow displayed source sum and zero-inclusive natural-number convention.
% BN-SRC-680 TARGET-CORRECTION-END
% BN-SRC-681 TARGET-CORRECTION-BEGIN
\par\smallskip\noindent\textnormal{\small\textbf{উৎস-সংশোধন (BN-SRC-681):} সূত্রের চলরাশির জায়গায় উৎসের অসম্পূর্ণ ভাষাগত !! চিহ্ন ছিল। ওই স্থানগুলিতে সূত্র-চলরাশি A পুনরুদ্ধার করা হয়েছে; নামযুক্ত ভাষাগত টোকেনগুলি অপরিবর্তিত।} % BN-SRC-681 TARGET-CORRECTION-END
% BN-SRC-682 TARGET-CORRECTION-BEGIN
\par\smallskip\noindent\textnormal{\small\textbf{উৎসসংশোধন (BN-SRC-682):} উৎসে সব সূত্র-অপারেটরকে দুটি সূত্র যুক্তকারী বলা হয়েছে; নঞর্থক ও পরিমাণক অপারেটর একটিমাত্র সূত্রের উপর কাজ করে। গদ্যে অপারেটরের নিজস্ব স্থানসংখ্যা রাখা হয়েছে।}
% BN-SRC-682 TARGET-CORRECTION-END
% BN-NORM-182 TARGET-NORMALIZATION-BEGIN
% Bengali display annotations retained; optional normalization note withdrawn.
% BN-NORM-182 TARGET-NORMALIZATION-END

\end{document}
